% Appendix -- control documents and the three measurement traps.
\section{Control documents in full}
\label{app:controls}

$\dvec$ is defined by the document it is subtracted against, so ``content'' is
whatever separates the skill from that control. Our default control holds only
length fixed. Two tighter ones, generated from the skill document itself, hold
more: \textbf{shuffled} keeps every word and every number and destroys the line
order; \textbf{corrupted} keeps the wording, the procedure, the table layout and
the ordering of the units, and replaces \emph{every conversion factor}. All
three are injected under the same name and header, so the prompt never announces
which it is.

Both tighter controls move the accounting a long way
(Table~\ref{tab:controls}). A table whose numbers are all wrong solves $0.720$
of the rescued cell; a table with the right numbers in scrambled order solves
$0.659$; the unrelated text solves $0.098$. Against no document the cell is
$0.000$ by construction and with the gold document it is $1.000$. So when either
tight control is used, $\gvec$ is no longer inert---it repairs $0.63$--$0.74$---
and the residual $\dvec$ repairs only $0.17$--$0.28$.

\begin{table*}[tbp]
\caption{The same decomposition against three control documents. Synthetic
control, multiple choice, Qwen3-1.7B, $n=358$, rescued cell $n_R=82$, every
one of layers 21--27. ``in context'' is the control document actually placed in the prompt;
``$+\gvec$'' is its state injected at the final position of the no-document
prompt. That the two agree to within $0.03$ for three very different documents
is the strongest validity check we have on the patch itself. $\dvec$ is the
same measurement in all three rows; only what it is subtracted against changes.}
\label{tab:controls}
\centering
\small
\begin{tabular}{lccc}
\toprule
control document & \multicolumn{3}{c}{fraction of the rescued cell correct}\\
\cmidrule(lr){2-4}
 & control in context & $h_{\text{no}}+\gvec$ & $h_{\text{no}}+\dvec$ \\
\midrule
unrelated text, matched length      & 0.098 & \textbf{0.098} & \textbf{0.500--0.805} \\
same document, lines scrambled      & 0.659 & 0.634--0.671 & 0.171--0.280 \\
same document, factors replaced     & 0.720 & 0.720--0.744 & 0.171--0.256 \\
\midrule
gold skill in context & 1.000 & --- & --- \\
no document           & 0.000 & --- & --- \\
\bottomrule
\end{tabular}

\end{table*}

\paragraph{The same control on third-party skills.} We repeated the corrupted
control on the clinical calculators, where the arm already exists in the
released harness (numeric constants in the prose perturbed, wording, structure
and executable tools untouched). That material is free-form and answered through
a chain of thought, and there the corrupted skill is worth \emph{nothing}: on 20
rescued items it scores $0.050$, which is what no skill at all scores, against
$0.950$ for the gold skill. The synthetic and the real material therefore agree,
and they agree on the thing that matters---whether corrupting a document costs
you anything is a property of the answer format, not of the document.

Read as an attribution rather than as three competing estimates, this says
something uncomfortable about the multiple-choice form of the task. On the full
358 items the gold document is worth $0.179 \to 0.288$; a scrambled version of
it is worth $0.179 \to 0.268$ and a version with every factor wrong is worth
$0.179 \to 0.237$. Four fifths of the measured skill effect survives destroying
the order of the document, and just over half survives replacing every number
in it. What the model needs, in this format, is mostly a document containing the
right vocabulary and the right kinds of number---not a correct table. The
correct table is worth the remaining $0.02$--$0.05$ absolute, or $0.17$--$0.28$
of the rescued cell.

That is a statement about the answer format as much as about the model. We
registered the test before running it: in the free-form format the number has to
be produced, so a corrupted table should be worth nothing. It is
(Table~\ref{tab:format}). Corrupting every factor keeps $0.54$ of the
multiple-choice effect on the full item set and \emph{none} of the free-form
effect's $0.467$. Scrambling the line order costs little in either, which is what it
should: our shuffle permutes lines, and each row of the conversion table is one
line, so every factor is still there to be retrieved---just not in a readable
order.

\subsection{Multiple choice manufactures a floor and merges retrieval with
execution}
\label{sec:format}

Our synthetic items exist in two forms with identical content: four-option
multiple choice, and free-form (``give the number only''). The forms do not
measure the same thing, and there is a mechanism in the literature for why.
\citet{wiegreffe2025mcq} localise the choice of an answer \emph{symbol} to a
single middle layer and one to four attention heads per layer: in multiple
choice the decision is assembled into one position and one token. That is
exactly the regime in which a single-position patch transports the whole task,
and it is why our final-position channel reads $1.000$ there and $0.000$ under a
chain of thought on the same items.



Inspecting what the small model writes explains the top half. Without the
document it echoes the quantity from the question (``2''); with the document it
writes a \emph{conversion factor from the table} (``21'', ``180'', ``60''). It
retrieves and does not execute. Multiple choice scores retrieval as success
whenever the retrieved number can be matched against an option, and the
positional prior supplies a floor of $0.18$--$0.29$ on top.

The bottom half is the sharper version, and it is a prediction we registered
before running: if the multiple-choice form is solvable from a document's
structure rather than its numbers, a table with every factor replaced must score
near zero in free-form. It scores exactly zero.

Two consequences. A benchmark comparing skill effects across model scales in
multiple-choice format is measuring a mixture whose composition changes with
scale. And an ablation that perturbs a document and reports the accuracy drop is
measuring, in multiple choice, mostly whether the document is still recognisably
a document.

\section{Three measurement traps, in full}
\label{app:traps}

\subsection{The aggregate gold-answer log-probability ranks the controls first}

Averaged over all items, the log-probability of the gold answer under each
patched condition at layer 27, the last, orders as
\begin{center}
mean vector $-2.85$ \; $>$ \; filler state $-5.20$ \; $>$ \; \textbf{true skill
state} $-6.60$ \; $>$ \; no document $-7.30$ \; $>$ \; wrong skill $-10.29$,
\end{center}
i.e.\ a content-free average vector scores $3.75$ nats \emph{better} than the
state the model actually builds while reading the document, with accuracy
$0.256$ against that state's $0.436$.

The cells of \S\ref{sec:cells} account for the gap arithmetically
(Table~\ref{tab:cells}): F and B---where the document does not work---contribute
$133\%$ of it at layer 21, where the gap is $2.75$ nats, while on R the true state wins by $2.5$ nats and $0.69$ in
accuracy. The mechanism is visible in the option distribution: the mean vector
raises the entropy over the four answer tokens from $0.187$ to $0.972$ nats
(maximum $\ln 4 = 1.386$). A flattened distribution cannot assign any token
$-19$ nats, so it looks excellent exactly where the true state is confidently
wrong, and flattening never promotes a token to the argmax, so accuracy does not
move. The one point the two channels agree on is that the \emph{wrong skill} is
worst on both---positive evidence for content specificity, since injecting the
wrong content actively misleads.

The practical rule is that an aggregate log-probability over a mixed item set
measures confidence calibration as much as it measures content, and that a
decomposition by outcome cell is the cheapest way to tell which.

\subsection{A band where the forward pass is damaged}

Between layers 17 and 20 every patched condition improves together and the
conditions become indistinguishable. Reading by cell isolates the cause: on the
K cell---items the model already answers correctly---the true-state patch drops
the gold log-probability from $-1.6$ at layer 16 to $-7.1$ at layer 18, while
the filler-state patch, built over an equally long prompt, stays at $-1.5$. Any
transfer or recovery number falling in 17--20 is reported with that band marked.

\subsection{What the wrong-document control is actually worth}
\label{sec:plausible}


Holding the prompt's length, structure and the document slot's position fixed and
changing only what sits in that slot gives an ordered comparison the usual
binary---document or no document---cannot make. On the rescued items, with
everything else identical ($n=40$, intervals in brackets): no document at all
$0.025$ $[0.000, 0.073]$; non-clinical prose in the slot $0.250$
$[0.116, 0.384]$; this skill's own paired neutral document $0.175$
$[0.057, 0.293]$; one fixed wrong clinical skill $0.150$ $[0.039, 0.261]$; the
gold document $0.950$ $[0.882, 1.000]$.

Two things follow, one solid and one not. The solid one is that the
wrong-document control is not the same object as the no-document control: any
document in the slot moves these items off the floor, so a paper that reports
``with the document'' against ``without'' is measuring presence and content
together, and the presence part is not negligible on a cell selected for failing
without. That is the reason every arm in \S\ref{sec:battery} is measured against
a receiver that holds a document rather than against an empty slot.

The one that is not solid is the ordering. A clinical distractor costing more
than irrelevant prose ($0.150$ against $0.250$) is the direction we expected---a
near miss should be worse than an obvious miss, and the three conditions fall in
that order---but every pair of intervals overlaps at $n=40$, so we report a
direction and not a result. We flag it because the design that would settle it is cheap and
this material affords it: these 55 skills contain graded near misses by
construction (four different formulae for the corrected QT interval, two for the
delta gap), so ``how similar must a wrong document be before it hurts'' can be
measured on an ordered scale rather than a binary. We have not run that
experiment.



